% !TEX root = ../main.tex
\chapter{モナド：失敗・状態・非同期を合成する形}
\label{chap:monad}

\begin{chapterabstract}
本章では、モナドを「文脈を持つ処理を合成する形」として導入する。
モナドを抽象的な数学用語として暗記するのではなく、失敗する処理、エラー理由を返す処理、状態を読み書きする処理、非同期に完了する処理を、同じ合成パターンで扱うための道具として読む。
Lean 4 では、まず自作の小さな \texttt{Option} 用 \texttt{pure} と \texttt{bind} を定義し、モナド則を等式として確認する。
最後に、失敗しうるユーザー移行を題材に、\texttt{bind} の法則が effectful code のリファクタリング規則として働くことを見る。
\end{chapterabstract}

\begin{learninggoals}
\item \texttt{Option}、\texttt{Except}、状態つき計算を「文脈を持つ値」として説明できる。
\item \texttt{pure} と \texttt{bind} の役割を、実務コードの処理連結として読める。
\item モナド則を、抽象法則ではなく effectful code のリファクタリング規則として説明できる。
\item Lean 4 で \texttt{Option} の小さなモナド則を証明できる。
\item 失敗しうるマイグレーションを \texttt{bind} で安全に合成する最小モデルを書ける。
\end{learninggoals}

\section{この章を読む理由}

実務コードでは、普通の関数だけでは表しにくい処理が頻繁に出てくる。
たとえば、入力が欠けていれば失敗するバリデーション、失敗理由を返すマイグレーション、内部状態を更新する処理、外部 API の応答を待つ非同期処理である。
これらは見た目には別物に見える。
しかし、どれも「値だけでなく、値が置かれている文脈も一緒に運ぶ処理」として見ると、同じ形で合成できる。

\begin{intuitionbox}
普通の関数 \texttt{A -> B} は、\texttt{A} の値を受け取って \texttt{B} の値を返す。
一方、失敗しうる関数 \texttt{A -> Option B} は、\texttt{B} の値を返すかもしれないし、返せないかもしれない。
この「かもしれない」という文脈を保ったまま次の処理につなぐための道具が \texttt{bind} である。
\end{intuitionbox}

\FigurePlaceholder{fig-ch17-overview.png}{0.90}{モナドを文脈つき処理の合成として見る}{fig:ch17-overview}

この章で扱うモナドは、Mathlib の高度な圏論 API としてのモナドではない。
本章では、初学者がコードとして追えるように、\texttt{Option}、\texttt{Except}、状態つき計算の最小モデルだけを使う。
圏論的には、モナドは自己関手と二つの自然変換を持つ構造として定義される。
ただし本書の実務入門では、まず \texttt{pure} と \texttt{bind} を使った合成規則として理解すれば十分である。

\section{\texttt{Option}、\texttt{Except}、\texttt{State} を同じ形で見る}

最初に、三つの典型例を並べる。
ここで大事なのは、細部ではなく「値をそのまま返すのではなく、何らかの文脈に入れて返す」という共通点である。

\begin{center}
\begin{tabular}{p{0.22\linewidth}p{0.30\linewidth}p{0.38\linewidth}}
\toprule
型の形 & 文脈 & 実務での読み方 \\
\midrule
\texttt{Option A} & 値があるかもしれない & 欠損、検索失敗、バリデーション失敗 \\
\texttt{Except E A} & 失敗理由 \texttt{E} を持つ & エラーコード、移行失敗理由、API エラー \\
\texttt{S -> (A, S)} & 状態 \texttt{S} を読み書きする & カウンタ、セッション、更新履歴、疑似的な状態遷移 \\
非同期結果 & 後で値が得られる & タスク、Future、Promise、外部 I/O \\
\bottomrule
\end{tabular}
\end{center}

\begin{softwarebox}
\texttt{Option} は「失敗したらそこで止まる」処理に向いている。
\texttt{Except} は「失敗した理由も残したい」処理に向いている。
状態つき計算は「値だけでなく、新しい状態も次へ渡したい」処理に向いている。
非同期処理は「値が今すぐではなく、後で得られる」処理として読める。
どれも、普通の関数合成 \texttt{g (f x)} だけでは扱いにくい文脈を持っている。
\end{softwarebox}

本章では、このような文脈を抽象的に \texttt{M A} と書くことにする。
\texttt{A} は中身の型であり、\texttt{M} は文脈を表す。
たとえば、\texttt{M A} が \texttt{Option A} なら「失敗するかもしれない \texttt{A}」である。
\texttt{M A} が \texttt{Except E A} なら「エラー \texttt{E} を出すかもしれない \texttt{A}」である。

\section{\texttt{pure} と \texttt{bind}}

モナドの操作として、まず \texttt{pure} と \texttt{bind} を見る。
\texttt{pure} は、普通の値を文脈の中に入れる。
\texttt{bind} は、文脈つきの値を取り出せた場合だけ、次の文脈つき処理へ渡す。

\begin{definitionbox}
本章では、モナドを次の二つの操作を持つものとして扱う。
\begin{itemize}
\item \texttt{pure : A -> M A}。普通の値を文脈 \texttt{M} に入れる。
\item \texttt{bind : M A -> (A -> M B) -> M B}。文脈つきの値を、次の文脈つき処理へつなぐ。
\end{itemize}
この二つが自然に振る舞うために、後で述べる三つのモナド則を要求する。
\end{definitionbox}

\texttt{Option} の場合、\texttt{pure} は単に \texttt{some} で値を包む。
\texttt{bind} は、入力が \texttt{none} なら次の処理を実行せずに \texttt{none} を返す。
入力が \texttt{some a} なら、\texttt{a} を次の処理へ渡す。

\begin{leanbox}
ここでは Lean 標準の \texttt{Option.bind} を直接使わず、理解のために \texttt{optPure} と \texttt{optBind} を自作する。
これにより、\texttt{bind} が何をしているかを \texttt{match} で直接確認できる。
\end{leanbox}

\begin{listing}[htbp]
\begin{minted}{lean4}
def optPure {α : Type} (a : α) : Option α :=
  some a

def optBind {α β : Type} (x : Option α) (f : α → Option β) : Option β :=
  match x with
  | none => none
  | some a => f a

#eval optBind (some 10) (fun n => some (n + 1))  -- => some 11
#eval optBind (none : Option Nat) (fun n => some (n + 1))  -- => none
\end{minted}
\caption{Option 用の pure と bind}
\label{lst:ch17_option_bind}
\end{listing}

コード \ref{lst:ch17_option_bind} の一つ目の \texttt{\#eval} は、\texttt{some 10} の中身を取り出して \texttt{some 11} を作る。
二つ目は、最初から \texttt{none} なので、後続の関数を実行せずに \texttt{none} になる。
この「失敗を自動的に伝播する」振る舞いが、バリデーションやマイグレーションで重要になる。

\FigurePlaceholder{fig-ch17-bind-flow.png}{0.90}{\texttt{bind} が成功時だけ次の処理へ進む流れ}{fig:ch17-bind-flow}

\section{モナド則をリファクタリング規則として読む}

モナドには三つの基本法則がある。
数学としては、これらは \texttt{pure} と \texttt{bind} が一貫して振る舞うことを表す。
ソフトウェア的には、これらは effectful code のリファクタリング規則である。
つまり、不要な包みを消してもよいこと、何もしない後続処理を消してもよいこと、括弧の付け方を変えてもよいことを保証する。

\begin{center}
\begin{tabular}{p{0.23\linewidth}p{0.34\linewidth}p{0.35\linewidth}}
\toprule
法則 & 形 & リファクタリングとしての読み方 \\
\midrule
左単位律 & \texttt{bind (pure a) f = f a} & 値を一度包んですぐ渡すだけなら、直接渡してよい \\
右単位律 & \texttt{bind x pure = x} & 最後にただ包み直すだけの処理は消してよい \\
結合律 & \texttt{bind (bind x f) g = bind x (fun a => bind (f a) g)} & 処理列の括弧の付け方を変えても意味は変わらない \\
\bottomrule
\end{tabular}
\end{center}

\begin{softwarebox}
モナド則は「実行順序を自由に入れ替えてよい」という意味ではない。
特に状態更新や非同期処理では、順序そのものは重要である。
モナド則が保証するのは、同じ順序の処理列について、\texttt{bind} の括弧や不要な \texttt{pure} を整理しても意味が変わらない、ということである。
\end{softwarebox}

\FigurePlaceholder{fig-ch17-monad-laws.png}{0.92}{モナド則を effectful code のリファクタリング規則として読む}{fig:ch17-monad-laws}

\section{Lean で \texttt{Option} の小さな性質を証明する}

\texttt{Option} の三つの法則は、場合分けで証明できる。
\texttt{x : Option A} には \texttt{none} と \texttt{some a} の二通りしかない。
Lean の \texttt{cases} は、この二通りを漏れなく分解する。

\begin{leanbox}
左単位律は定義を展開すればすぐ同じ式になるので \texttt{rfl} で証明できる。
右単位律と結合律は、\texttt{x} が \texttt{none} か \texttt{some} かで分けると、どちらのケースも \texttt{rfl} で閉じる。
\texttt{<;>} は、直前のタクティクで発生したすべてのゴールに同じ後続タクティクを適用する記法である。
\end{leanbox}

\begin{listing}[htbp]
\begin{minted}{lean4}
theorem option_left_identity {α β : Type} (a : α) (f : α → Option β) :
    optBind (optPure a) f = f a := by
  rfl

theorem option_right_identity {α : Type} (x : Option α) :
    optBind x optPure = x := by
  cases x <;> rfl

theorem option_assoc {α β γ : Type}
    (x : Option α) (f : α → Option β) (g : β → Option γ) :
    optBind (optBind x f) g =
      optBind x (fun a => optBind (f a) g) := by
  cases x <;> rfl
\end{minted}
\caption{Option のモナド則}
\label{lst:ch17_option_laws}
\end{listing}

ここで証明しているのは、\texttt{Option} の中身について特別な性質ではない。
どんな型 \texttt{A}、\texttt{B}、\texttt{C} であっても、\texttt{Option} の失敗伝播の形そのものがこの法則に従う、という主張である。
この点が、テストと証明の違いである。
特定の入力例だけではなく、\texttt{Option} の全パターンについて確認している。

\section{サンプル：失敗するマイグレーションを安全に合成する}

次に、失敗しうるマイグレーションを考える。
旧データ \texttt{UserDraft} には、メールアドレスと年齢があるかもしれないし、ないかもしれない。
新データ \texttt{UserV2} では、メールアドレスと年齢を必須にしたい。
この場合、移行は常に成功するとは限らない。

\begin{softwarebox}
実務では、これは「旧データに欠損があると新スキーマへ完全には移せない」という状況である。
ここで重要なのは、失敗ケースを例外的な裏口として扱わず、返り値の型 \texttt{Option UserV2} に明示することである。
型に失敗可能性を入れると、後続処理はその失敗可能性を無視できなくなる。
\end{softwarebox}

\begin{listing}[htbp]
\begin{minted}{lean4}
structure UserDraft where
  id : Nat
  emailOpt : Option String
  ageOpt : Option Nat
  deriving Repr, DecidableEq

structure UserWithEmail where
  id : Nat
  email : String
  ageOpt : Option Nat
  deriving Repr, DecidableEq

structure UserV2 where
  id : Nat
  email : String
  age : Nat
  deriving Repr, DecidableEq
\end{minted}
\caption{移行前後の小さなデータモデル}
\label{lst:ch17_user_model}
\end{listing}

移行を一気に書くのではなく、二段階に分ける。
まずメールアドレスを必須化し、次に年齢を必須化する。
どちらも失敗する可能性があるので、返り値は \texttt{Option} になる。

\begin{listing}[htbp]
\begin{minted}{lean4}
def requireEmail (u : UserDraft) : Option UserWithEmail :=
  match u.emailOpt with
  | none => none
  | some email => some { id := u.id, email := email, ageOpt := u.ageOpt }

def requireAge (u : UserWithEmail) : Option UserV2 :=
  match u.ageOpt with
  | none => none
  | some age => some { id := u.id, email := u.email, age := age }

def migrateUser (u : UserDraft) : Option UserV2 :=
  optBind (requireEmail u) requireAge
\end{minted}
\caption{失敗しうる移行ステップを bind で合成する}
\label{lst:ch17_user_migration}
\end{listing}

\texttt{requireEmail u} が \texttt{none} なら、\texttt{requireAge} は実行されない。
\texttt{requireEmail u} が \texttt{some withEmail} なら、\texttt{withEmail} が \texttt{requireAge} へ渡される。
これにより、失敗の検査を各ステップに閉じ込めながら、全体の移行を一本の処理として組み立てられる。

\begin{listing}[htbp]
\begin{minted}{lean4}
def goodDraft : UserDraft :=
  { id := 7, emailOpt := some "a@example.com", ageOpt := some 20 }

def missingEmailDraft : UserDraft :=
  { id := 8, emailOpt := none, ageOpt := some 20 }

def missingAgeDraft : UserDraft :=
  { id := 9, emailOpt := some "b@example.com", ageOpt := none }

#eval migrateUser goodDraft  -- => some { id := 7, email := "a@example.com", age := 20 }
#eval migrateUser missingEmailDraft  -- => none
#eval migrateUser missingAgeDraft  -- => none
\end{minted}
\caption{成功例と失敗例}
\label{lst:ch17_user_examples}
\end{listing}

この例では、成功時だけ \texttt{some UserV2} が得られる。
メールアドレスが欠けている場合も、年齢が欠けている場合も、結果は \texttt{none} になる。
\texttt{Option} では失敗理由は残らない。
失敗理由まで必要なら、次節の \texttt{Except} を使う。

\FigurePlaceholder{fig-ch17-migration-pipeline.png}{0.92}{失敗しうる移行ステップを \texttt{bind} で合成する}{fig:ch17-migration-pipeline}

\subsection{do 記法は bind の読みやすい書き方}

Lean では、\texttt{Option} や \texttt{Except} のようなモナド的な型に対して \texttt{do} 記法を使える。
これは \texttt{bind} を読みやすく書くための構文である。
次の二つの移行関数は、同じ意味を持つ。

\begin{listing}[htbp]
\begin{minted}{lean4}
def migrateUserDo (u : UserDraft) : Option UserV2 := do
  let withEmail ← requireEmail u
  requireAge withEmail

theorem migrateUser_refactor (u : UserDraft) :
    migrateUser u = migrateUserDo u := by
  cases u with
  | mk id emailOpt ageOpt =>
      cases emailOpt <;> cases ageOpt <;> rfl
\end{minted}
\caption{bind 版と do 記法版の同値性}
\label{lst:ch17_do_notation}
\end{listing}

\begin{syntaxbox}
\texttt{do} は、段階的な処理を上から順に書くための構文である。
\begin{itemize}
\item \texttt{x <- action} は、文脈つきの計算 \texttt{action} から中身を取り出して \texttt{x} として使う。
\item \texttt{Option} の文脈では、途中で \texttt{none} になると後続の処理へ進まない。
\item 本質的には \texttt{bind} を読みやすく書くための記法である。
\end{itemize}
\end{syntaxbox}


証明では、\texttt{UserDraft} の中身を分解し、\texttt{emailOpt} と \texttt{ageOpt} の有無を場合分けしている。
これは「すべての欠損パターンについて、\texttt{bind} 版と \texttt{do} 記法版が同じ結果になる」と確認している、という意味である。

\section{\texttt{Except}：失敗理由を残す}

\texttt{Option} は、失敗したかどうかだけを表す。
一方、実務では失敗理由が必要になることが多い。
欠損メールなのか、欠損年齢なのか、フォーマット不正なのかを区別したい。
その場合は \texttt{Except E A} を使う。

\begin{definitionbox}
\texttt{Except E A} は、成功時に \texttt{A} の値を返し、失敗時に \texttt{E} 型のエラーを返す型である。
\texttt{Option A} が「失敗するかもしれない \texttt{A}」なら、\texttt{Except E A} は「理由 \texttt{E} つきで失敗するかもしれない \texttt{A}」である。
\end{definitionbox}

\begin{listing}[htbp]
\begin{minted}{lean4}
inductive MigrationError where
  | missingEmail
  | missingAge
  deriving Repr, DecidableEq

def requireEmailE (u : UserDraft) : Except MigrationError UserWithEmail :=
  match u.emailOpt with
  | none => Except.error MigrationError.missingEmail
  | some email => Except.ok { id := u.id, email := email, ageOpt := u.ageOpt }

def requireAgeE (u : UserWithEmail) : Except MigrationError UserV2 :=
  match u.ageOpt with
  | none => Except.error MigrationError.missingAge
  | some age => Except.ok { id := u.id, email := u.email, age := age }
\end{minted}
\caption{Except による失敗理由つき移行}
\label{lst:ch17_except}
\end{listing}

\texttt{Except} でも、成功時だけ次へ進み、失敗時はそのエラーを保ったまま止まる。
\texttt{Option} と \texttt{Except} の違いは、失敗の情報量である。
合成の形は同じである。

\begin{listing}[htbp]
\begin{minted}{lean4}
def exceptBind {ε α β : Type}
    (x : Except ε α) (f : α → Except ε β) : Except ε β :=
  match x with
  | Except.error e => Except.error e
  | Except.ok a => f a

theorem except_right_identity {ε α : Type} (x : Except ε α) :
    exceptBind x exceptPure = x := by
  cases x <;> rfl

theorem except_assoc {ε α β γ : Type}
    (x : Except ε α) (f : α → Except ε β) (g : β → Except ε γ) :
    exceptBind (exceptBind x f) g =
      exceptBind x (fun a => exceptBind (f a) g) := by
  cases x <;> rfl
\end{minted}
\caption{Except 用の bind とモナド則}
\label{lst:ch17_except_laws}
\end{listing}

このコードでは、\texttt{Except.error}なら後続処理へ進まず、\texttt{Except.ok}なら中身を取り出して関数\texttt{f}へ渡す\texttt{exceptBind}を定義している。
\texttt{except\_right\_identity}と\texttt{except\_assoc}は、失敗理由を持つ計算でも\texttt{Option}と同じ形の合成法則が成り立つことを示している。

\section{\texttt{State}：値と一緒に状態を渡す}

状態つき計算も、同じ発想で扱える。
状態 \texttt{S} を受け取り、値 \texttt{A} と新しい状態 \texttt{S} を返す関数を考える。
これは \texttt{S -> (A, S)} という形になる。

\begin{definitionbox}
本章では、状態つき計算を \texttt{SimpleState S A = S -> (A, S)} として表す。
これは「状態 \texttt{S} を受け取り、値 \texttt{A} と更新後の状態 \texttt{S} を返す処理」である。
\end{definitionbox}

\begin{listing}[htbp]
\begin{minted}{lean4}
abbrev SimpleState (σ α : Type) := σ → (α × σ)

def statePure {σ α : Type} (a : α) : SimpleState σ α :=
  fun s => (a, s)

def stateBind {σ α β : Type}
    (m : SimpleState σ α) (f : α → SimpleState σ β) : SimpleState σ β :=
  fun s =>
    let r := m s
    f r.fst r.snd
\end{minted}
\caption{状態つき計算の最小モデル}
\label{lst:ch17_state}
\end{listing}

\texttt{stateBind} は、まず計算 \texttt{m} を現在の状態 \texttt{s} で実行する。
その結果として、値 \texttt{r.fst} と新しい状態 \texttt{r.snd} が得られる。
次に、値を後続処理 \texttt{f} へ渡し、更新後の状態でその処理を実行する。

\begin{listing}[htbp]
\begin{minted}{lean4}
def getState : SimpleState Nat Nat :=
  fun s => (s, s)

def putState (s : Nat) : SimpleState Nat Unit :=
  fun _ => ((), s)

def tick : SimpleState Nat Nat :=
  stateBind getState (fun n =>
    stateBind (putState (n + 1)) (fun _ =>
      statePure n))

#eval tick 10  -- => (10, 11)
\end{minted}
\caption{カウンタ状態を更新する例}
\label{lst:ch17_state_example}
\end{listing}

この例では、\texttt{tick 10} は現在の状態 \texttt{10} を値として返し、状態を \texttt{11} に更新する。
状態つき処理では、順序が特に重要である。
ただし、同じ順序の処理列について、\texttt{bind} の括弧を整理できることはリファクタリング上大きな意味を持つ。

\section{非同期処理をどう読むか}

非同期処理も、モナド的な合成として読める。
たとえば、\texttt{Task A} や \texttt{Future A} のような型を考えると、それは「今すぐではないが、将来 \texttt{A} が得られる計算」である。
\texttt{pure} は、すでに手元にある値を完了済みの非同期結果として包む。
\texttt{bind} は、前の非同期結果が完了してから、その値を使う次の非同期処理を開始する。

\begin{warningbox}
本章では、実際のスケジューリング、キャンセル、タイムアウト、並行実行、外部 I/O の意味論までは形式化しない。
非同期を扱う現実のシステムでは、順序、失敗、再試行、リソース解放などが別途重要になる。
ここでは、非同期処理も「文脈を持つ値を次の処理へ渡す」という形で読める、という範囲に留める。
\end{warningbox}

\section{実務上の読み替え：合成則をレビューの言葉にする}

モナドを実務に持ち込むとき、最初から「これはモナドである」と言う必要はない。
むしろ、次の問いに答えられることが重要である。

\begin{itemize}
\item この処理の返り値には、どのような文脈が付いているか。
\item 失敗、エラー、状態、非同期性を、返り値の型に明示しているか。
\item 次の処理へ進む条件は \texttt{bind} の形で一箇所に閉じ込められているか。
\item 不要な \texttt{pure} や括弧の整理をしても、意味が変わらない形になっているか。
\item 失敗理由が必要なのに \texttt{Option} だけで潰していないか。
\end{itemize}

\begin{softwarebox}
例外処理、バリデーション、状態更新、非同期 API 呼び出しは、どれも「途中で普通の関数合成から外れる」処理である。
モナド的な見方を使うと、これらを場当たり的な制御フローではなく、合成可能な処理列として設計できる。
\end{softwarebox}

\section{よくある誤解}

\begin{warningbox}
モナドは「副作用をなくす魔法」ではない。
失敗、状態、非同期性を消すのではなく、それらを型と合成規則の中に明示する道具である。
また、モナド則は処理順序を勝手に入れ替えてよいという許可ではない。
同じ順序の処理列について、\texttt{pure} と \texttt{bind} の構文的な整理が安全であることを述べている。
\end{warningbox}

\section{本章のまとめ}

\texttt{Option A}、\texttt{Except E A}、\texttt{S -> (A, S)} は、いずれも文脈を持つ値として読める。

\texttt{pure} は普通の値を文脈に入れ、\texttt{bind} は文脈つきの値を次の文脈つき処理へつなぐ。

モナド則は、effectful code の左単位、右単位、結合律というリファクタリング規則として読める。

Lean では、\texttt{Option} の \texttt{none} と \texttt{some} を場合分けすることで、モナド則を小さく証明できる。

失敗しうるマイグレーションでは、失敗可能性を返り値の型に入れ、\texttt{bind} で段階的に合成すると、失敗ケースを漏らしにくくなる。

失敗理由が必要な場合は \texttt{Option} ではなく \texttt{Except} を使う。

状態つき処理や非同期処理も、値だけでなく文脈を次へ渡す合成として読むことができる。
